\clearpage
\properlane{hearing-king}{The King Makes a People Capable of Hearing Him}
The kingdom Jesus describes to Pilate is present among human beings, yet its life comes from another source. Its subjects are made anew. The Mass expresses this change as rescue from darkness, healing of division, obedient hearing and perseverance toward an enduring inheritance. Augustine's exposition of John explains entry into this kingdom; Jerome's exposition of Daniel explains the lasting life its members receive in the incarnate Son.

\subsection{In the world, but with another origin}
Pilate hears a claim to kingship within the practical categories of his office. A king has subjects, interests and means of resistance. A claimant delivered by his own authorities and lacking an armed rescue appears either defeated or dangerous. Jesus' answer changes the question from the kingdom's military capacity to its origin. Augustine attends closely to the difference between a kingdom that is not \emph{in} the world and one that is not \emph{of} it. Jesus says the latter. The distinction allows Augustine to identify the kingdom with believers who remain in the world, while refusing to derive their new life from the world.\footnote{Augustine, \emph{Tractates on John} 115.1--2, NPNF I.7.}

This answer is concrete. A people who remain here still have neighbors, work, households, rulers and responsibilities. Their baptism does not remove them from human history. Augustine recalls Christ's prayer that his disciples be kept from evil while remaining in the world. He also recalls the mixture of wheat and tares until the harvest. The kingdom's present existence does not imply the present purification of every member or the disappearance of all evil within its visible reach. Its origin and its final perfection answer different questions.

The contrast between worldly origin and presence in the world is especially useful beside the Gradual. Its worldwide reach could otherwise be mistaken for a map of successful Christian administration. Augustine instead asks what makes the people a kingdom at all. They were under the prince of this world and become Christ's people through regeneration. At precisely this point he cites Colossians 1:13, the transfer into the beloved Son's kingdom. The connection between the Gospel and Epistle is thus already present in his actual argument: changed citizenship names a saving change in persons, not a relocation to another part of the earth.

The Collect gives that argument a communal voice. Sin is a wound that disperses human families. The petition for Christ's gentle dominion concerns the cure of that wound. A society can share an external authority while remaining divided by envy, fear and hatred. Augustine's account explains why the gathering asked for in the prayer reaches more deeply than an administrative settlement. The King creates willing allegiance through grace. Its gentleness does not make the claim less demanding; it identifies the healer whose demand restores the persons who obey him.

The Gospel's absence of armed rescue now has positive content. The disciples cannot bring this kingdom into existence by fighting to prevent Jesus' delivery. Such a rescue would oppose the very self-giving through which the Lamb makes a people. The Introit therefore supplies the proper beginning for this interpretation. The assembly praises someone who has been slain, and approaches him as the source of its life. Its first act is thanksgiving for a victory it did not manufacture.

\subsection{Why hearing is more than receiving information}
At the end of the appointed Gospel, Jesus connects belonging to truth with hearing his voice. The order of the statement troubles an easy account of conversion. One might suppose that a hearer independently recognizes a true proposition and thereby earns the status of belonging to truth. Augustine reads the sentence in the opposite direction. A person hears because the truth has given him a share in itself. Hearing is the fruit of grace, not a prior achievement that obliges grace to follow.\footnote{Augustine, \emph{Tractates on John} 115.4.}

He reaches this conclusion by distinguishing two ways of being related to truth. Every human being is a creature made through the divine Word. That universal dependence cannot by itself explain the hearing named here, because not everyone believes. Jesus is speaking of a further gift. Augustine describes its effect as the inward hearing of the person who obeys and believes. The hearer has not merely received sound or understood the grammar of an assertion. The heart is turned toward the speaker as trustworthy Lord.

The distinction protects both grace and responsibility. If hearing meant only exposure to words, Pilate's conversation would already constitute full obedience. If it meant an independent achievement preceding grace, salvation would become the reward of a merit the sinner supplied unaided. Augustine's account instead identifies an enabled response. The listener really believes and obeys, but that response begins in Christ's gracious calling. Thanksgiving, which opens the Epistle, is consequently the fitting speech of someone who has become a hearer.

Chrysostom's reading of Colossians 1:12 makes the same gift vivid from another angle: the Father has not simply offered an inheritance but made the recipients fit to share it. A gift can stand before someone who has no capacity to use it; divine generosity also changes the recipient. The transfer from darkness is therefore richer than release from a prison. It brings the rescued into a royal relationship.\footnote{Chrysostom, \emph{Homilies on Colossians} II, exposition of Col 1:12--15.}

This is why the Preface's truth, life, holiness and grace belong together. Truth is not detached information about a remote ruler. It is the truth of the Son who gives life and sanctifies those he gathers. Holiness describes the new relation's character; grace identifies its gratuitous source. The qualities of the kingdom answer the Gospel's question about the sort of King standing before Pilate, while the Gospel gives those qualities a personal center.

A concrete moral consequence follows. To say that one belongs to Christ while refusing every correction is to divide allegiance from hearing. To treat another person's conversion as a trophy of one's own cleverness is to forget the prior gift. The Mass trains a different response: gratitude for having been made fit, patient witness to others, and willingness to have one's own habits corrected. These are consequences of Augustine's account of hearing, not conditions by which a sinner purchases the first grace.

\subsection{How the eternal Son receives dominion}
Daniel's vision presents another difficulty. The human figure receives power and a kingdom before the Ancient of Days. If Christ is eternally divine, why is dominion given to him? Jerome addresses the problem by attending to the figure's humanity. The Son of Man signifies the Son of God's assumption of human flesh. The giving of power, honor and kingdom is read through the Incarnation and the self-emptying described by Paul in Philippians 2.\footnote{Jerome, \emph{Commentary on Daniel}, VII, on 7:13--14, PL 25, cols. 670--671.}

The divine Son has not advanced from the condition of a creature into divinity. He has truly assumed the condition in which obedience, humiliation and exaltation can be predicated of him. Jerome's reasoning therefore refuses to solve the difficulty by making the humanity unreal. The Son's received glory belongs to the saving economy of the one who became man. The very humanity through which he suffers is not abandoned when he reigns.

This distinction matters for the people Augustine describes. Their hope is not merely to watch a divine power remain secure above their vulnerability. The incarnate King brings human nature into the enduring kingdom. Daniel 7:18 and 27 associate the saints with its possession. Jerome contrasts their heavenly inheritance with kingdoms that arise from earth and return to it. He rejects an account that exhausts the promise in Maccabean victory, because the promised universality and permanence exceed that event; he also rejects an earthly millennial kingdom as the final form of the saints' inheritance.\footnote{Jerome, \emph{Commentary on Daniel}, VII, on 7:17--18 and 27.}

Augustine and Jerome thus answer different parts of one question. Augustine explains how former subjects of darkness become Christ's people. Jerome explains why their kingdom can be received through a genuinely human King without making his divine nature inferior or their final inheritance another temporary empire. The Alleluia's endurance gives their changed membership a future. The Postcommunion can ask for shared heavenly reign because the Son has made humanity his own and the saints belong to his kingdom.

The Introit's acclaim of the Lamb can be heard within the same doctrinal relation. The slain one receives honor; the Epistle identifies him as uncreated source of creatures. His redemptive victory makes divine glory manifest in the humanity he has assumed. Receiving royal acclaim and possessing eternal divinity are not competing descriptions of two different Christs. The priestly King of the Preface is one person, and his people live from his one saving work.

\subsection{Intercession gathers what sin has scattered}
The Offertory turns from the assembly's hearing to the Son's asking. The Father promises the nations as his inheritance. Augustine reads Psalm 2:8 within the temporal economy of incarnate sacrifice and intercession. The Son asks for the salvation of those who will be his possession. His people are not assembled by the self-assertion of one nation against another, but through the priestly agency of their King.\footnote{Augustine, \emph{Expositions of the Psalms}, Ps 2.7.}

This makes the Offertory a fitting passage from proclamation into sacrifice. The peoples named by the Collect and Gradual enter the horizon of petition. Their number and distance do not make them strangers to the Son's concern. Nor can the faithful congratulate themselves as though they already represented a humanity fully reconciled. The Secret asks for unity and peace precisely because the gifts remain necessary. The one offered is himself the giver who can gather them.

Bellarmine's treatment of Psalm 2 helps locate the struggle involved. The Lord's authority can break the habits that enslave a person. Correction is mercy when it destroys a sinful attachment and frees the person to obey; judgment also remains real for resistance that refuses the King. The image cannot be transferred into a permission for the congregation to break other people. The Gospel has supplied the decisive image of royal action: Christ gives himself, and his servants hear.\footnote{Bellarmine, \emph{Commentary on Psalms}, Ps 2, comments on vv. 7--10.}

\subsection{The peace of pilgrims}
The Communion's blessing could be misunderstood as a promise that obedience will produce an untroubled life. Psalm 28 has prepared a stronger account. The voice of the Lord resounds through disruption, and the King remains. Augustine reads the concluding verses as the Church's preservation amid the world's flood and its future reign with Christ. Peace in Christ is compatible with adversity in the world.\footnote{Augustine, \emph{Expositions of the Psalms}, Ps 28.10--11.}

The difference changes how perseverance is understood. A difficult week does not prove that the blessing failed. It may be precisely the setting in which the received gift must be lived. Bellarmine likewise distinguishes the strength needed on the journey from the final blessing toward which it leads. The Postcommunion places nourishment before continued service and shared reign. Those who have heard the King's voice need continuing sustenance; their first response does not make further dependence unnecessary.

The movement of this reading is therefore from redemption to regenerated hearing, from hearing to intercession for all peoples, and from sacramental nourishment to faithful perseverance. Every royal claim concerns a King who makes his subjects capable of life with him. Universal reach keeps the argument from shrinking into private feeling; the heavenly inheritance keeps it from identifying salvation with a successful political settlement.

\begin{fourSenses}
\item[Literal.] Jesus answers a royal accusation before Pilate; Colossians gives thanks for deliverance and transfer; Daniel presents an enduring kingdom; the prayers ask for healed allegiance, peace and perseverance.
\item[Allegorical.] The incarnate Son gathers the Church through redemption and regeneration. His received human glory opens the enduring kingdom to the saints whose life comes from him.
\item[Moral.] Hearing becomes obedient faith. Gratitude displaces self-sufficiency, intercession reaches other peoples, and correction of sinful habits serves faithful life together.
\item[Anagogical.] The pilgrim Church receives peace amid present conflict and hopes for lasting reign with Christ. Its final inheritance surpasses every temporary earthly kingdom.
\end{fourSenses}
